

\textbf{Robustness to Model Incompleteness and Deviations.}
To evaluate the system's resilience to discrepancies between the \textit{as-planned} BIM and the \textit{as-built} site, we simulate two scenarios: model incompleteness and plan deviations. First, we simulate partially built environments by omitting 20\% and 30\% of the BIM walls from the matching process. Second, we perturb all BIM walls\addtwo{, including the two used for initial alignment, to evaluate the system’s tolerance to \textit{as-built} vs. \textit{as-planned} discrepancies and the robustness of the initial alignment}. For each wall, the specific deviation is sampled from a uniform distribution, \changetwo{$[0, \Delta d]$ and $[0, \Delta \theta]$}{$U[-\Delta d, +\Delta d]$ for position and $U[-\Delta \theta, +\Delta \theta]$ for orientation}. \addtwo{To avoid dependence on a single perturbation sample, we repeat each sequence and condition on five deviated BIM models, each generated by re-sampling the per-wall deviations.
}

\textbf{Runtime Analysis.}
To verify real-time operation, we measured the processing frame rate (FPS) and the time required for initial alignment across four representative sequences.


